% 一時お試し: エネルギー収支を成分明示で書き下すとどう見えるか確認するだけ
\documentclass[12pt]{article}
\usepackage{luatexja}
\usepackage[match]{luatexja-fontspec}
\setmainjfont{HaranoAjiGothic}
\usepackage{amsmath}
\usepackage{graphicx}
\usepackage[a4paper,margin=12mm]{geometry}
\pagestyle{empty}
\begin{document}

\textbf{(1) 現行：ベクトル表記}
\[
\frac{dT}{dz}=-\frac{(1-\varepsilon)\left(\tfrac{\pi}{4}D^{2}\right)\boldsymbol{\Delta H}^{\top}\boldsymbol{r}}{\boldsymbol{F}^{\top}\boldsymbol{C}_p}
\]

\vspace{1.5em}
\textbf{(2) 成分明示・1行（横に長い）}
\[
\resizebox{\textwidth}{!}{$\displaystyle
\frac{dT}{dz}=-\frac{(1-\varepsilon)\left(\tfrac{\pi}{4}D^{2}\right)\left(\Delta H_1 r_1+\Delta H_2 r_2+\Delta H_3 r_3\right)}
{F_{\mathrm{C_3H_8}}C_{p,\mathrm{C_3H_8}}+F_{\mathrm{C_3H_6}}C_{p,\mathrm{C_3H_6}}+F_{\mathrm{H_2}}C_{p,\mathrm{H_2}}+F_{\mathrm{C_2H_4}}C_{p,\mathrm{C_2H_4}}+F_{\mathrm{CH_4}}C_{p,\mathrm{CH_4}}+F_{\mathrm{C_2H_6}}C_{p,\mathrm{C_2H_6}}}
$}
\]

\vspace{1.5em}
\textbf{(3) 成分明示・分母を $\sum$ で（コンパクト）}
\[
\frac{dT}{dz}=-\frac{(1-\varepsilon)\left(\tfrac{\pi}{4}D^{2}\right)\left(\Delta H_1 r_1+\Delta H_2 r_2+\Delta H_3 r_3\right)}
{\displaystyle\sum_{i} F_i\,C_{p,i}},\qquad i\in\{\mathrm{C_3H_8,\,C_3H_6,\,H_2,\,C_2H_4,\,CH_4,\,C_2H_6}\}
\]

\vspace{1.5em}
\textbf{(4) 分子も反応名で（参考）}
\[
\resizebox{\textwidth}{!}{$\displaystyle
\boldsymbol{\Delta H}^{\top}\boldsymbol{r}=
\underbrace{\Delta H_1 r_1}_{\text{脱水素}}+\underbrace{\Delta H_2 r_2}_{\text{クラッキング}}+\underbrace{\Delta H_3 r_3}_{\text{水素化}}
$}
\]

\end{document}
